\documentclass[aps,prb,twocolumn,10pt,raggedbottom,nofootinbib,superscriptaddress]{revtex4-2}
\usepackage{amsmath,amssymb,amsthm,mathtools,bm}
\usepackage{booktabs}
\usepackage[colorlinks=true,linkcolor=blue,citecolor=blue,urlcolor=blue]{hyperref}
\usepackage{microtype}
\hypersetup{pdftitle={Relaxation gap of a trajectory-averaged fermionic reset circuit},
  pdfauthor={Class-A fermionic adaptive-circuit project}}
\newcommand{\Id}{\hat{\mathbf{1}}}
\newcommand{\id}{\mathbf{1}}
\newcommand{\Tr}{\operatorname{Tr}}
\newcommand{\E}{\widetilde{\mathcal E}}
\begin{document}
\title{Relaxation gap of a trajectory-averaged fermionic reset circuit}
\author{Class-A fermionic adaptive-circuit project}
\noaffiliation
\date{October 1, 2026}
\begin{abstract}
We relate the relaxation of the outcome-averaged density matrix to the
single-particle projector product $A$ for one measurement cycle. The
dimensionless channel gap is $g=1-r_{\rm MB}$, distinct from the
logarithmic decay rate $\Delta=-\log r_{\rm MB}$. For perfect occupation
resets in a fixed repeated schedule, the charge-neutral and parity-even
many-body gaps equal $g_C=1-\rho(A)^2$. Higher-order correlations cannot
relax more slowly than the two-point function. The result allows
noncommuting, support-truncated Wannier measurements.
\end{abstract}
\maketitle
\setcounter{tocdepth}{1}
\tableofcontents

\section{Measurement cycle and two-point relaxation}
After discarding outcomes and fresh ancillary modes, the physical-layer
channel fills lower-band overcomplete Wannier (OW) modes and empties
upper-band modes, as in Bhuiyan, Pan, and Jian (BPJ), Appendix~D~\cite{BPJ}.
We retain exact discrete resets in their prescribed order.

Let $\{\hat\psi_i,\hat\psi_j^\dagger\}=\delta_{ij}$ for $N$ fermion
modes. At step $j$, measure $\hat\chi_j=\sum_i v_{ji}\hat\psi_i$,
with $v_j^\dagger v_j=1$ and
$\hat N_j=\hat\chi_j^\dagger\hat\chi_j$. Averaging the measurement
and perfect feedback gives the physical-layer channels
\begin{align}
 \E_{j,0}(\hat\rho)&=\hat\chi_j\hat\rho\hat\chi_j^\dagger
 +(\Id-\hat N_j)\hat\rho(\Id-\hat N_j),\nonumber\\
 \E_{j,1}(\hat\rho)&=\hat\chi_j^\dagger\hat\rho\hat\chi_j
 +\hat N_j\hat\rho\hat N_j. \label{eq:reset}
\end{align}
The target occupation is $s_j=0$ or $1$. The projector terms include
measurement dephasing; no separate dephasing step is added. One cycle
is the composition $\E=\E_{M,s_M}\circ\cdots\circ\E_{1,s_1}$,
with step~1 first. Distinct OW modes need not be orthogonal, so this
order matters.

Write $\overline G_{ab}=\Tr(\overline{\hat\rho}\,
\hat\psi_a^\dagger\hat\psi_b)$ and set
$P_j=v_jv_j^\dagger$, $Q_j=\id-P_j$. A reset removes correlations
with the measured mode and sets its occupation:
\begin{equation}
 \overline G'=Q_j\overline GQ_j+s_jP_j. \label{eq:localG}
\end{equation}
The part of the initial two-point function surviving one cycle is
therefore propagated by
\begin{equation}
 A=Q_M\cdots Q_1,\qquad
 \overline G_{n+1}=A\overline G_nA^\dagger+B. \label{eq:cycle}
\end{equation}
The source $B$ depends on the reset targets, not on the input state.
For deviations from a stationary correlation matrix,
$\delta G_{n+1}=A\delta G_nA^\dagger$. Column vectorization gives
$A^*\otimes A$, whose eigenvalues are $a_i a_j^*$ if $a_i$ are those
of $A$. Define the spectral radius $r=\rho(A)=\max_i|a_i|$.
No eigenvalue of $A$ is excluded: $A$ is not the many-body channel,
and stationarity is the affine condition $G_{\rm ss}=AG_{\rm ss}A^\dagger+B$,
not an eigenvalue 1 of $A$. The two-sided action accounts for the square:
the largest covariance multiplier is $r^2$. For $0<r<1$,
\begin{equation}
 g_C=1-r^2,\qquad \Delta_C=-2\log r. \label{eq:rate}
\end{equation}
Here $g_C$ is dimensionless and $\Delta_C$ is per complete cycle.
The spectral radius, not the largest singular value, sets asymptotic
decay. In the computation we diagonalize both invariant hard-wall blocks
of the ordered product and take the largest eigenvalue modulus across
both. A unit-modulus eigenvalue would give $g_C=0$; none is discarded.

\section{From two-point decay to the many-body gap}
Fix an invariant operator sector containing the stationary state.
For a mixing channel, $\E(\hat\rho_{\rm ss})=\hat\rho_{\rm ss}$
provides a simple eigenvalue 1. A nonstationary eigenoperator satisfies
$\E(\hat X_\nu)=\mu_\nu\hat X_\nu$, so its amplitude after $n$
cycles is $\mu_\nu^n$. Define
\begin{align}
 r_{\rm MB}&=\max_{\nu\text{ nonstationary}}|\mu_\nu|,
 &g_{\rm MB}&=1-r_{\rm MB},\nonumber\\
 \Delta_{\rm MB}&=-\log r_{\rm MB},
 &g_{\rm MB}&=1-e^{-\Delta_{\rm MB}}. \label{eq:definition}
\end{align}
The absolute channel gap is the deficit from unity in \emph{modulus},
not $\min_\nu|1-\mu_\nu|$ for complex eigenvalues. Only the stationary
mode is removed; another unit-modulus mode would obstruct mixing.
The gap and decay rate agree only to leading order near zero.

Covariance closure alone gives a bound. The Heisenberg channel
preserves the span of $\Id$ and all
$\hat\psi_a^\dagger\hat\psi_b$, so its restriction has eigenvalues
$1$ and $a_i a_j^*$. These also belong to the full channel spectrum.
In any mixing sector containing these operators,
\begin{equation}
 r_{\rm MB}\ge r^2,\quad g_{\rm MB}\le g_C,\quad
 \Delta_{\rm MB}\le\Delta_C.
 \label{eq:bound}
\end{equation}
This is the spectral-inclusion bound: a positive two-point rate alone
cannot exclude a slower higher-order correlation.

For $A^\dagger u=a^*u$, the centered occupation
\begin{equation}
 \hat O_u=\sum_{ab}u_a^*u_b
 \bigl(\hat\psi_a^\dagger\hat\psi_b-(G_{\rm ss})_{ab}\Id\bigr)
 \label{eq:densitymode}
\end{equation}
obeys $\E^\dagger(\hat O_u)=|a|^2\hat O_u$. Choosing $|a|=r$
explicitly realizes the slow neutral density mode.

Choose a basis starting with $\hat\chi_j$ and let $\hat B$ contain
only orthogonal modes. For even total operators, Eq.~\eqref{eq:reset} gives
\begin{align}
 \E_{j,s}^\dagger(\hat B)&=\hat B &&(\hat B\text{ even}),\nonumber\\
 \E_{j,s}^\dagger(\hat\chi_j^{(\dagger)}\hat B)&=0
 &&(\hat B\text{ odd}),\nonumber\\
 \E_{j,s}^\dagger(\hat N_j\hat B)&=s\hat B
 &&(\hat B\text{ even}). \label{eq:localrule}
\end{align}
A reset erases a measured leg; an occupation pair contributes $s$ and
reduces the degree by two. Thus $2p$-point functions couple only to
themselves and lower orders, without assuming Wick factorization.

Charge-neutral correlations take the form
\begin{equation}
 F_{I;J}^{(p)}=\Tr\!\left(\hat\rho\,
 \hat\psi_{i_1}^\dagger\cdots\hat\psi_{i_p}^\dagger
 \hat\psi_{j_p}\cdots\hat\psi_{j_1}\right), \label{eq:moments}
\end{equation}
where $I,J$ are increasing $p$-element sets. The order-$2p$ part
propagates creation indices with $A$ and annihilation indices with
$A^*$. Antisymmetry gives the one-cycle block
\begin{equation}
 H_p[I,J;K,L]=\det A[I,K]\,\det A[J,L]^*. \label{eq:Hp}
\end{equation}
Cauchy--Binet composes the minors of the $Q_j$ in the actual order.
The hierarchy is block triangular in $p$; lower-order terms do not
change its eigenvalues. The multipliers are
\begin{equation}
 \Lambda_{I,J}=\prod_{i\in I}a_i\prod_{j\in J}a_j^*,\qquad
 |I|=|J|=p. \label{eq:products}
\end{equation}
The independent sets may overlap, allowing the same label on both
legs. These $\sum_p\binom Np^2=\binom{2N}{N}$ moments span all neutral
operators. Schur triangularization gives the same products for
nondiagonalizable $A$.

The constant ($p=0$) supplies eigenvalue 1. Every $p\ge1$ product
obeys $|\Lambda_{I,J}|\le r^{2p}\le r^2$, with equality at $p=1$ for
$I=J=\{i_*\}$, $|a_{i_*}|=r$. No higher neutral correlation is slower;
the stationary eigenvalue is simple and
\begin{align}
 \boxed{g_{q=0}=g_C=1-\rho(A)^2,} \label{eq:neutralgap}\\
 \Delta_{q=0}=\Delta_C=-2\log\rho(A). \nonumber
\end{align}

\section{Sectors and earlier results}
Charge neutral means $[\hat N_F,\hat\rho]=0$, where
$\hat N_F=\sum_i\hat\psi_i^\dagger\hat\psi_i$; it permits mixtures
of number sectors. Definite-charge reset Kraus operators preserve this
condition while changing particle number. Maximal mixing and
fixed-number Slater preparations both lie in this sector.

With $p$ creations and $\ell$ annihilations, the even-sector argument
gives products of $p+\ell$ factors. Every nonconstant even product
has at least two, while the neutral sector attains $r^2$. Hence
$g_{\rm even}=g_{q=0}$ and $\Delta_{\rm even}=\Delta_{q=0}$.
To include odd operators, use the
similarity $\mathcal S(\hat O)=(-1)^{\hat N_F}\hat O$. It changes the
sign of the odd-Kraus sandwich in
$\mathcal S\E_j^\dagger\mathcal S$; on odd operators that sign cancels
the fermionic reordering sign and restores Eq.~\eqref{eq:localrule}.
Odd products then have maximal modulus $r$, attained with one factor.
For the specified Kraus extension on unrestricted operator space,
\begin{equation}
 g_{\rm full}=1-r,\qquad \Delta_{\rm full}=-\log r=\Delta_C/2.
\end{equation}
Physical density matrices are parity even. Thus the plotted $g_C$ is
the relevant many-body channel gap for our states, not merely a bound.
Earlier plots of $\Delta_C$ show the corresponding decay rate instead;
neither quantity is the occupation-spectrum gap around $1/2$.

The continuous-time counterpart is Barthel--Zhang, Sec.~III.6,
Proposition~7 and Eq.~(63): the covariance generator is a block of the
many-body Liouvillian~\cite{BZ}. Their Sec.~III.5, Proposition~6 and
discussion after Eq.~(57d), distinguish unrestricted and two-point
rates. Equation~\eqref{eq:bound} is the discrete spectral-inclusion
bound (also Ref.~\cite{project}, Eq.~(4)); equality~\eqref{eq:neutralgap}
requires our additional exact-reset argument.

BPJ Appendix~D, Eqs.~(D34)--(D40), instead treats effective
continuous-time bulk dissipation rates~\cite{BPJ}, not this ordered-reset
identity. Nonnormal transients and Jordan prefactors remain possible:
our result fixes an asymptotic rate, not a size-uniform mixing or
trajectory-purification time. It does not prove a thermodynamic
extrapolation. Extra dephasing, imperfect feedback, or randomized
schedules requires a new higher-order analysis.

\begin{thebibliography}{3}
\bibitem{BPJ}
A. Bhuiyan, H. Pan, and C.-M. Jian,
\href{https://doi.org/10.1103/nk38-ygyq}{Phys. Rev. Research \textbf{8}, 023147 (2026)},
Appendix D.
\bibitem{BZ}
T. Barthel and Y. Zhang,
\href{https://doi.org/10.1088/1742-5468/ac8e5c}{J. Stat. Mech. (2022) 113101},
Secs. III.5--III.6;
\href{https://arxiv.org/pdf/2112.08344}{arXiv:2112.08344v5}, pp. 14--16.
\bibitem{project}
A. Bhuiyan, \textit{A spectral bound on relaxation of the
trajectory-averaged circuit}, project note (September 29, 2026), Eq. (4);
\href{../../../matched_markov_lindblad_campaign/spectral_gap_v1/results/20260929T191448Z/analysis/gap_note.pdf}{local PDF}.
\end{thebibliography}
\end{document}
